% Festlegung des Allgemeinen Dokumentenformats
\documentclass[a4paper,12pt]{article}

% Umlaute unter UTF8 nutzen
\usepackage[utf8]{inputenc}

% weitere Pakete
% Grafiken aus PNG Dateien einbinden
\usepackage{graphicx}

% Deutsche Sonderzeichen und Silbentrennung nutzen
\usepackage[ngerman]{babel}

% Eurozeichen einbinden
\usepackage[right]{eurosym}

% Zeichenencoding
\usepackage[utf8]{inputenc}

%\usepackage{lmodern}

% floatende Bilder ermöglichen
%\usepackage{floatflt}

% mehrseitige Tabellen ermöglichen
\usepackage{longtable}

% Unterstützung für Schriftarten
%\newcommand{\changefont}[3]{ 
%\fontfamily{#1} \fontseries{#2} \fontshape{#3} \selectfont}

% Paket für Boxen im Text
\usepackage{fancybox}

% bricht lange URLs "schön" um
\usepackage[hyphens,obeyspaces,spaces]{url}

% Paket für Textfarben
\usepackage{color}

% Mathematische Symbole importieren
\usepackage{amssymb}

% auf jeder Seite eine Überschrift (alt, zentriert)
%\pagestyle{headings}

% erzeugt Inhaltsverzeichnis mit Querverweisen zu den Abschnitten (PDF Version)
%\usepackage[bookmarksnumbered,pdftitle={\titleDocument},hyperfootnotes=false]{hyperref}
%\hypersetup{colorlinks, citecolor=red, linkcolor=blue, urlcolor=black}
%\hypersetup{colorlinks, citecolor=black, linkcolor= black, urlcolor=black}

% für Tabellen
\usepackage{array}
\usepackage{amsmath}

% Runde Klammern für Zitate
%\usepackage[numbers,round]{natbib}

% Paket für Zeilenabstand
\usepackage{setspace}

% für Bildbezeichner
\usepackage{capt-of}

% für Listings
\usepackage{listings}

\usepackage[scaled]{helvet}
\usepackage{etoolbox}

%tocloft für weitere möglichkeiten beim TOC (zb bold sections)
%\usepackage{tocloft}

% Variablen
%Variablen welche innerhalb der gesamten Arbeit zur Verfügung stehen sollen
\newcommand{\titleDocument}{Hausarbeit_JonasTröger}
\newcommand{\subjectDocument}{im Fach Software Engineering}


% Packet für Seitenrandabstände und Einstellung für Seitenränder
\usepackage{geometry}
\geometry{left=3.0cm, right=2.0cm, top=2.5cm, bottom=2.0cm}

\usepackage{fancyhdr} %Paket laden
% neue Kopfzeilen mit fancypaket
\usepackage{lastpage}
\pagestyle{fancy} %eigener Seitenstil
\fancyhf{} %alle Kopf- und Fußzeilenfelder bereinigen

%\fancyhead[L]{\nouppercase{\leftmark}} %Kopfzeile links
\fancyhead[C]{} %zentrierte Kopfzeile
\fancyhead[R]{Seite \thepage \hspace{1pt} von \pageref{LastPage} } %Kopfzeile rechts

\renewcommand{\headrulewidth}{0.4pt} %obere Trennlinie
%\fancyfoot[C]{\thepage} %Seitennummer
%\renewcommand{\footrulewidth}{0.4pt} %untere Trennlinie

% Festlegung Art der Zitierung - Havardmethode: Abkuerzung Autor + Jahr
\bibliographystyle{plainurl}

% Schaltet den zusätzlichen Zwischenraum ab, den LaTeX normalerweise nach einem Satzzeichen einfügt.
\frenchspacing

%Konfiguriere das Inhaltsverzeichnis
\usepackage{tocbasic}
\DeclareTOCStyleEntries[raggedentrytext,numwidth=0pt,numsep=1ex,dynnumwidth,]
{tocline}{chapter,section,subsection,subsubsection,paragraph,subparagraph}
\DeclareTOCStyleEntries[
  indent=0pt,
  linefill=\TOCLineLeaderFill,
]{tocline}{section,subsection,subsubsection,paragraph,subparagraph}

\usepackage{tocloft}

\setlength{\cftsecindent}{0pt}
\setlength{\cftsubsecindent}{0pt}
\setlength{\cftsubsubsecindent}{0pt}
\setcounter{tocdepth}{3}         								% Level Depth
\renewcommand{\cftsecfont}{\bfseries}				%Typewriter sectionentries
%\renewcommand{\cftsubsecfont}{\ttfamily \itshape}

\lstset{numbers=left, numberstyle=\tiny, numbersep=5pt, keywordstyle=\color{black}\bfseries, stringstyle=\ttfamily,showstringspaces=false,basicstyle=\footnotesize,captionpos=b}
\lstset{language=java}

% für Stichwortverzeichnis
\usepackage{makeidx}
% Indexerstellung
\makeindex

%Kein indent nach linebreak
\setlength{\parindent}{0pt}

\renewcommand\familydefault{\sfdefault}

\patchcmd{\thebibliography}{*}{}{}{}

\graphicspath{ {./input/abb/} }

%/////////////////////////////////////////////////////////////////////////////////////////////////////////////////////////////////////
\begin{document}

\renewcommand{\figurename}{Abb.}
% hier werden die Trennvorschläge inkludiert
%%hier müssen alle Wörter rein, welche Latex von sich auch nicht korrekt trennt bzw. bei denen man die genaue Trennung vorgeben möchte
\hyphenation{
Film-pro-du-zen-ten
Lux-em-burg
Soft-ware-bau-steins
zeit-in-ten-siv
}


% Schriftart Helvetica verwenden
%\usepackage{helvet}
%\renewcommand\familydefault{\sfdefault}

% Leere Seite am Anfang
%\thispagestyle{empty} % erzeugt Seite ohne Kopf- / Fusszeile
%\mbox{}
%\newpage


% Titelseite %
\thispagestyle{empty}
\onehalfspacing

\begin{verbatim}


\end{verbatim}

\begin{center}
\begin{large}
%\textbf{Jonas Tröger}\\
%\vspace{6mm}
-Universität-\\
-Standort-
\end{large}
\end{center}

\begin{verbatim}




\end{verbatim}
\begin{center}
\textbf{\Large{Thema}}\\
\vspace{6mm}
\large{Bachelorarbeit \linebreak 
zur Erlangung der staatlichen Abschlussbezeichnung eines \linebreak 
Bachelor of Science (B. Sc.)\linebreak
 in der Studienrichtung Informatik}
\end{center}

\vfill

\begin{flushleft}
\noindent \begin{tabular}{ll}
Eingereicht von: & -Vorname Nachname- \\ 
 & -Adresse- \\ 
 & -Seminargruppe- \\ 
 & -Student nr- \\ 
 Betreuer: & -Titel- -Vorname Nachname- \\
 & -Adresse (Uni oder Arbeit)- \\
 
\end{tabular}
\end{flushleft}
\vspace{10mm}
\noindent  -Ort-, den \today\\



% römische Numerierung
%\pagenumbering{latin}

% 1.5 facher Zeilenabstand
\onehalfspacing

%\newpage

% Sperrvermerk
%\thispagestyle{empty}
%\input{sperrvermerk}

%\newpage
% Einleitung / Abstract
%\thispagestyle{empty}
%\input{0_abstract}

% arabische Seitennummerierung ab hier
\pagenumbering{arabic}

% einfacher Zeilenabstand
%\singlespacing

\newpage
% Seitenzählung bei Inhaltsverzeichnis beginnen
\setcounter{page}{1}

% Inhaltsverzeichnis anzeigen
\thispagestyle{empty}
\tableofcontents


%%%%%%% EINLEITUNG %%%%%%%%%%%%
%\newpage
%\fancyhead[L]{\nouppercase{\leftmark}} %Kopfzeile links

% 1,5 facher Zeilenabstand
\onehalfspacing

% Alternative Einbindung des Abstract in Kapitel "0" falls gewünscht
%\setcounter{section}{-1}
%\setcounter{page}{0}

% Option: Einbindung abstract
%\input{0_abstract}
%\newpage

% einzelne Kapitel werden hier eingebunden
%\input{Kapitel/Geschichte}
\newpage
\section{kapitel1}
texttextexttexttextexttexttextexttexttextexttexttextexttexttext
exttexttextexttexttextexttexttextexttexttextexttexttextext
texttextexttexttextexttexttextexttexttextexttexttextexttexttexte
ttexttextexttexttextexttexttextexttexttextexttexttextext
texttextexttexttextexttexttextexttexttextexttexttextexttexttexte
xttexttextexttexttextexttexttextexttexttextexttexttextext
\subsection{kapitel1.1}
texttextexttexttextexttexttextexttexttextexttexttextexttexttex
texttexttextexttexttextexttexttextexttexttextexttexttextext
texttextexttexttextexttexttextexttexttextexttexttextexttextte
xtexttexttextexttexttextexttexttextexttexttextexttexttextext
\subsubsection{kapitel1.1.1}
texttextexttexttextexttexttextexttexttextexttexttextexttextte
xtexttexttextexttexttextexttexttextexttexttextexttexttextext
texttextexttexttextexttexttextexttexttextexttexttextexttextte
xtexttexttextexttexttextexttexttextexttexttextexttexttextext
texttextexttexttextexttexttextexttexttextexttexttextexttextte
xtexttexttextexttexttextexttexttextexttexttextexttexttextexts


\newpage
\section{Kapitel2}
texttextexttexttextexttexttextexttexttextextt
exttextexttexttextexttexttextexttexttextexttexttextexttexttextexttexttextext
texttextexttexttextexttexttextexttexttextexttext
textexttexttextexttexttextexttexttextexttexttextexttexttextexttexttextext
texttextexttexttextexttexttextexttexttextexttextt
extexttexttextexttexttextexttexttextexttexttextexttexttextexttexttextext
texttextexttexttextexttexttextexttexttextexttextte
xtexttexttextexttexttextexttexttextexttexttextexttexttextexttexttextext
\subsection{kapitel2.1}
texttextexttexttextexttexttextexttexttextexttexttext
exttexttextexttexttextexttexttextexttexttextexttexttextexttexttextext
texttextexttexttextexttexttextexttexttextexttexttexte
xttexttextexttexttextexttexttextexttexttextexttexttextexttexttextext
texttextexttexttextexttexttextexttexttextexttexttextex
ttexttextexttexttextexttexttextexttexttextexttexttextexttexttextext
texttextexttexttextexttexttextexttexttextexttexttextext
texttextexttexttextexttexttextexttexttextexttexttextexttexttextext

\subsubsection{kapitel2.1.1}
texttextexttexttextexttexttextexttexttextexttexttextex
ttexttextexttexttextexttexttextexttexttextexttexttextexttexttextext
texttextexttexttextexttexttextexttexttextexttexttextex
ttexttextexttexttextexttexttextexttexttextexttexttextexttexttextext
texttextexttexttextexttexttextexttexttextexttexttextext
texttextexttexttextexttexttextexttexttextexttexttextexttexttextext



%/////////////////////////////////////////////////////////////////////////////////////////////////////////////////////////////////////

%Zum einbinden ins Quellenverzeichnis nachfolgendes für alle Zitate / Quellen
%\nocite{bild1}


% einfacher Zeilenabstand
\singlespacing
% Literaturliste soll im Inhaltsverzeichnis auftauchen
\newpage
%\addcontentsline{toc}{section}{Literaturverzeichnis / Quellenverzeichnis}
% Literaturverzeichnis anzeigen
\renewcommand\refname{Quellenverzeichnis}
\bibliography{Bachelorarbeit.bib}

%% Index soll Stichwortverzeichnis heissen
% \newpage
% % Stichwortverzeichnis soll im Inhaltsverzeichnis auftauchen
% \addcontentsline{toc}{section}{Stichwortverzeichnis}
% \renewcommand{\indexname}{Stichwortverzeichnis}
% % Stichwortverzeichnis endgültig anzeigen
% \printindex

%\onehalfspacing
% evtl. Anhang
%\newpage
%\phantomsection
%\addcontentsline{toc}{section}{Anhang}
%\fancyhead[L]{Anhang} %Kopfzeile links
%\input{anhang/anhang}

% Eidesstattliche Erklärung
%\newpage
%\addcontentsline{toc}{section}{Selbständigkeitserklärung}
\newpage

\section{Selbstständigkeitserklärung}

Ich versichere, dass ich die vorliegende Arbeit ohne fremde Hilfe
selbstständig verfasst und nur die angegebenen Quellen und
Hilfsmittel benutzt habe. Wörtlich oder dem Sinn nach aus anderen
Werken entnommene Stellen sind unter Angabe der Quellen kenntlich
gemacht. Die Arbeit wurde bisher in gleicher oder ähnlicher Form
weder veröffentlicht, noch einer anderen Prüfungsbehörde vorgelegt.
\vfill
-Vorname Nachname-, -Ort-, \today \hfill 
\begin{tabular}{l}
\includegraphics[width=0.3\linewidth]{us}\\
\rule{120pt}{1pt}\\
\end{tabular}







% leere Abschlussseite
%\newpage
%\thispagestyle{empty} % erzeugt Seite ohne Kopf- / Fusszeile
%\mbox{}

\end{document}
